\documentclass[11pt, a4paper, titlepage, oneside]{scrbook}

\usepackage{../Tarc/Tarc}
\usepackage{../Tarc/Tarc-private}
\graphicspath{{immagini/}}

\title{Sistemi Operativi}
\author{Lorenzo Ferraro}
\date{\today}
\setDegree{Informatica}

\begin{document}

\frontmatter
\SapienzaCoverItalianoTriennale
\InfoSegnalazioni

\tableofcontents

\mainmatter

\chapter{Introduzione}

\section{Sistemi numerici}

Il sistema numerico usato normalmente è il sistema numerico \textbf{decimale} (o in base 10). In questo, ogni cifra rappresenta una potenza di \textbf{10}: $5374_{10} = 5 \cdot 10^3 + 3 \cdot 10^2 + 7 \cdot 10^1 + 4 \cdot 10^0$. Esistono, oltre a questo sistema, anche molti altri sistemi in basi diverse, tra cui il sistema \textbf{binario} e il sistema \textbf{esadecimale}.

\begin{Definizione}[Sistema binario]
\label{def-sistema-binario}
Il sistema \textbf{binario} rappresenta ogni numero tramite solo due cifre: \textbf{0} e \textbf{1}. Ogni cifra rappresenta una potenza di \textbf{2}: $1101_2 = 1 \cdot 2^3 + 1 \cdot 2^2 + 0 \cdot 2^1 + 1 \cdot 2^0 = 8 + 4 + 1 = 13$. 
\end{Definizione}

\begin{Definizione}[Sistema esadecimale]
\label{def-sistema-esadecimale}
Il sistema \textbf{esadecimale} rappresenta ogni numero tramite tutte le \textbf{10 cifre} e tramite le lettere che vanno dalla \textbf{A} alla \textbf{F} (che rappresentano i numeri da 10 a 15). Ogni cifra rappresenta una potenza di \textbf{16}: $4AF_{16} = 15 \cdot 16^0 + 10 \cdot 16^1 + 4 \cdot 16^2 =  15 + 160 + 1024 = 1199$.
\end{Definizione}

\subsection{Convertire i numeri tra i diversi sistemi numerici}

Abbiamo visto come convertire qualsiasi numero in binario/esadecimale in decimale, ma vediamo come fare il contrario.

\begin{itemize}
    \item \textbf{Decimale} $\longrightarrow$ \textbf{Binario}: vi sono due metodi:
    \begin{enumerate}[label=\arabic*°]
        \item Metodo: prendere la più grande potenza di 2 più piccola del numero da voler convertire, e sottrarla, e continuare così finché non si arriva a 0. Ad esempio, $68 = 64 - 4$, dunque $68_{10} = 1000100_2$.
        \item Metodo: dividere il numero per 2 e posizionare il resto appena ottenuto a sinistra della rappresentazione binaria. Ad esempio, $68/2 = 34 + 0 \longrightarrow 34/2 =  17 + 0 \longrightarrow 17/2 = 8 + 1 \longrightarrow 8/2 = 4 + 0 \longrightarrow 4/2 = 2 + 0 \longrightarrow 2/2 = 1 + 0 \longrightarrow 1/2 = 0 + 1$, dunque $68_{10} = 1000100_2$.
    \end{enumerate}
    \item \textbf{Decimale} $\longrightarrow$ \textbf{Esadecimale}: dividere  il numero per 16 e posizionare il resto appena ottenuto a sinistra della rappresentazione esadecimale. Ad esempio, $68/16 = 4 + 4 \longrightarrow 4/16 = 0 + 4$, dunque $68_{10} = 44_{16}$.
\end{itemize}

\begin{Proprieta}[Potenze di 2]
\label{prop-potenze-2}
Conoscendo solo alcune potenze di 2, è possibile approssimare i valori di molte più potenze. Infatti, se si conoscono tutte le potenze da $2^0$ a $2^9$, e sapendo che \mbox{$2^{10} \approx 1000, 2^{20} \approx 1000000, 2^{30} \approx 1000000000, \ldots$}, si possono approssimare tutte le potenze tra $2^0$ e $2^{39}$. Infatti, $2^{34} = 2^{30} \cdot 2^4 \approx 1000000000 \cdot 16 = 16000000000$.
\end{Proprieta}

\begin{Definizione}[Definizioni sistema binario]
\label{def-definizioni-sistema-binario}
Nel sistema binario, chiamiamo \textbf{bit} le cifre dei numeri. Un \textbf{nibble}, poi, sono 4 bit mentre un \textbf{byte} sono 8 bit. Esistono poi il \textbf{most significant bit} e \textbf{least significant bit}, che rispettivamente il bit più a sinistra e più a destra, e il \textbf{most significant byte} e \textbf{least significant byte}, che sono rispettivamente il byte più a sinistra e più a destra.  
\end{Definizione}

Ad esempio, nel numero $1010101010$, il primo 1 è il MSB, mentre l'ultimo 0 è il LSB. Allo stesso modo, nel numero CEB48E8 (rappresentato in esadecimale, dove ogni cifra rappresenta un nibble), il byte CE è il MSB, mentre il byte E8 è il LSB.

\subsection{Numeri col segno}

In qualsiasi sistema numerico, per rappresentare un numero negativo basta mettere il simbolo $-$ prima del numero. Questo, però, per i computer non è possibile: lavorando solo con dei bit, non si può anteporre ad un numero il simbolo $-$. Infatti, il computer non ha idea di cosa sia quel simbolo, conosce solo 0 e 1. Dunque bisogna trovare un'altra soluzione per rappresentare i numeri col segno. Vi sono due modi possibili:
\begin{itemize}
    \item \textbf{Segno e Ampiezza}
    \item \textbf{Completamento a 2}
\end{itemize}

\begin{Proprieta}[Segno e Ampiezza]
\label{def-Segno e Ampiezza}
\textbf{Segno e Ampiezza} è una tecnica usata dai computer per rappresentare i numeri col segno. Si usa il \textbf{bit più significativo per rappresentare il segno}: se questo vale \textbf{0}, allora il numero è \textbf{positivo}, altrimenti è \textbf{negativo}. La formula per convertire un numero in decimale è la seguente:
\[
\bm{A = (-1)^{a_{N-1}} \sum_{i = 0}^{N-2} a_i 2^i}
\]
\end{Proprieta}
Ad esempio, il numero 6 tramite il Segno e Ampiezza diventa in binario 0110, mentre -6 diventa 1110. È dunque necessario utilizzare un bit in più del normale per rappresentare i numeri. Con N bit si possono quindi rappresentare tutti i numeri da $-(2^{N-1} - 1)$ a $2^{N-1} - 1$. 

Il problema più grande di questa tecnica è l'addizione: semplicemente, questa non funziona. Se infatti proviamo a sommare -6 (1110) e 6 (0110), troviamo che
\[
\begin{array}{rccccccc}
& & & & & & \\
& & 1 & 1 & 1 & 0 & + &\\
& & 0 & 1 & 1 & 0 & = &\\
\hline
& 1 & 0 & 1 & 0 & 0& = & 19_{10} \neq 0
\end{array}
\]
Inoltre, un'altro problema è che vi sono due rappresentazioni dello 0: infatti, 0 = 1000 = 0000. 

Per questi motivi nasce un'altra tecnica, quella del \textbf{complemento a 2}.

\begin{Proprieta}[Complemento a 2]
\label{prop-complemento a 2}
\textbf{Complemento a 2} è un'altra tecnica usata per rappresentare i numeri col sengo. Il \textbf{bit più significativo} ha lo stesso valore visto con i numeri senza segno, ma con il \textbf{segno negativo}. La formula per convertire un numero in decimale è la segunete:
\[
\bm{A = a_{N-1}(-2^{N-1}) + \sum_{i = 0}^{N-2} a_i2^i}
\]
\end{Proprieta}

Ad esempio, il numero 7 tramite complemento a 2 diventa in binario 0111, mentre -7 diventa 1001. Con N bit si possono quindi rappresentare tutti i numeri da $-2^{N-1}$ a $2^{N-1} - 1$. 

Con questa implementazione si risolvono entrambi i problemi del Segno e Ampiezza. Infatti, lo zero è rappresentato solo dalla sequenza 0000, e l'addizione non ha problemi: infatti, se proviamo a sommare -6 (1010) e 6 (0110), troviamo che
\[
\begin{array}{rccccccc}
& & & & & & \\
& & 1 & 0 & 1 & 0 & + &\\
& & 0 & 1 & 1 & 0 & = &\\
\hline
& \cancel{1} & 0 & 0 & 0 & 0& = & 0
\end{array}
\]

dove l'uno inziale non si conta dato che stiamo parlando di numeri a 4 bit.

\begin{Proprieta}[Convertire un numero in un numero negativo tramite il Complemento a 2]
\label{prop-trasformare-complemento-2}
Preso un numero in binario possiamo facilmente trasformarlo nell'opposto tramite il Complemento a 2. Infatti, basta \textbf{invertire i bit} e \textbf{aggiungere 1}. Preso ad esempio $6_{10} = 0110_{2}$, troviamo che $-6_{10} = 1001_2 + 1 = 1010 = -6_{10}$.
\end{Proprieta}

Questa proprietà funziona per un motivo molto semplice: sappiamo che $6 + \bar{6} = 0110 + 1001 = 1111$, dunque aggiungendo un 1 troviamo che $6 + \bar{6} + 1 = 1111 + 1 = 10000$, e dato che il primo zero si elude, troviamo che $6 + \bar{6} + 1 = 0$, ossia $-6 = \bar{6} + 1 \longrightarrow \bm{-1001 = 0110 + 1}$.  

\section{Operazioni elementari}

\subsection{Addizioni e sottrazioni}

Le addizioni in binario e in esadecimale funzionano esattamente uguale a come funzionano in decimale. Ad esempio,

\noindent
\begin{minipage}{0.48\textwidth}
\[
\begin{array}{rcccc}
& \scriptstyle 1 & & & \\
& 0 & 1 & 1 & + \\
& 0 & 1 & 0 & = \\
\hline
& 1 & 0 & 1 &
\end{array}
\]
\end{minipage}\hfill
\begin{minipage}{0.48\textwidth}
\[
\begin{array}{rcccc}
& & \scriptstyle 1 & & \\
& A & 4 & F & + \\
& 2 & A & 1 & = \\
\hline
& C & F & 0 &
\end{array}
\]
\end{minipage}

Anche per la sottrazione il metodo rimane lo stesso. Ad esempio,

\noindent
\begin{minipage}{0.48\textwidth}
\[
\begin{array}{rcccc}
& & & \scriptstyle 1 (2) & \\
& 0 & \cancel{1} \scriptstyle 0 & 0 & - \\
& 0 & 0 & 1 & = \\
\hline
& 0 & 0 & 1 &
\end{array}
\]
\end{minipage}\hfill
\begin{minipage}{0.48\textwidth}
\[
\begin{array}{rcccc}
& & & \scriptstyle 1 (16) & \\
& 3 & \cancel{D} \scriptstyle C & 9 & - \\
& 1 & A & D & = \\
\hline
& 2 & 2 & C &
\end{array}
\]
\end{minipage}

\subsection{Il problema dell'overflow}

\begin{Osservazione}[Overflow]
\label{oss-overflow}
Un problema classico nelle addizioni tra numeri binari in un sistema digitale è l'\textbf{overflow}: i computer lavorano su un numero fisso di bit, e con un addizione è possibile sforare questo numero. 
\end{Osservazione}

Ad esempio, se lavoriamo con numeri a 4 bit, la somma $1001 + 1000 = 10001$ sfora, creando diversi problemi. Infatti, con N bit abbiamo un range di numeri rappresentabili, da $0$ a $2^N-1$ (questo se parliamo di numeri senza segno).\\
Vi sono due modi per risolvere questo problema:
\begin{itemize}
    \item \textbf{Sign-Extension}
    \item \textbf{Zero-Extension}
\end{itemize}

\begin{Proprieta}[Sign-Extension]
\label{prop-sign-extention}
La \textbf{sign-extension} è un'operazione usata dai computer per aumentare il numero di bit di un numero (in Complemento a 2) così da poter completare somme su di esso che superino il limite di bit prestabilito. La regola è semplice: \textbf{duplicare il bit più significativo a sinistra raddoppiando il numero di bit del numero} (raddoppiamo il numero di bit così da esser coperti sia sulle addizioni ma anche sulle moltiplicazioni, che vedremo tra poco).
\end{Proprieta}

Ad esempio, $3_{10} = 0011 = 00000011$, e $-5_{10} = 1011 = 1111011$. Questo con gli 0 ovviamente funziona, ma anche con gli 1: infatti, $-5 = 1011 = -8 + 2 + 1 = - 16 + 8 + 2 + 1 = -32 + 16 + 8 + 2 +1 = -64 + 32 + 16 + 8 +1 = 11111011$. 

\begin{Proprieta}[Zero-Extension]
\label{zero-extension}
La \textbf{zero-extension} è un'operazione usata dai coputer per aumentare il numero di bit di un numero così da poter completare somme su di esso che superino il limtie di bit prestabilito. La regola è semplice: \textbf{raddoppiare la lunghezza del numero aggiungendo tutti 0 a sinistra}.
\end{Proprieta}
Ad esempio, $3_{10} = 0011 = 00000011$ (come nella Sign-Extension). La Zero Extension però \textbf{non funziona con i numeri col segno}. Infatti, $-5_{10} = 1011 \neq 00001011 = 11_{10}$. 

\subsection{Moltiplicazioni}

Anche le moltiplicazioni in binario e in esadecimale funzionano come in decimale. Ad esempio,

\noindent
\begin{minipage}{0.48\textwidth}
\[
\begin{array}{rccccc}
& & 0 & 1 & 1 & \times \\
& & 0 & 0 & 1 & = \\
\hline
& & 0 & 1 & 1 & + \\
& 0 & 0 & 0 & & + \\
0 & 0 & 0 & & & + \\
\hline
0 & 0 & 0 & 1 & 1
\end{array}
\]
\end{minipage}\hfill
\begin{minipage}{0.48\textwidth}
\[
\begin{array}{rccccc}
& & 3 & B & 7 & - \\
& & 2 & C & 4 & = \\
\hline
& & E & D & C & + \\
& 2 & C & 9 & & + \\
7 & 6 & E & & & + \\
\hline
A & 4 & 6 & 1 & C
\end{array}
\]
\end{minipage}

\subsection{Shift di bit}

\begin{Definizione}[Shift]
\label{def-shift}
Uno \textbf{shift} è un'operazione sui bit che \textbf{sposta tutte le posizioni dei bit} di una sequenza binaria o a destra o a sinsitra di un determinato numero di posizioni.
\end{Definizione} 

Esistono due tipi di shift:
\begin{itemize}
    \item \textbf{shift aritmetico}, dove si ripete il bit più significativo sequendo la logica della sign-extension;
    \item \textbf{shift logico}, dove si inseriscono degli 0 a sinistra della sequenza sequendo la logica della zero-extension.
\end{itemize}

Ad esempio, $-7_{10} = 1001 \bm{>>} 2 = 0010 = 2_{10}$ è uno \textbf{shift logico} (si è infatti perso il segno, cosa che capita quando si usa la zero-extension), mentre $-8_{10} = 111000 \bm{>>>} 3 = 111111 = -1_{10}$ è uno \textbf{shift aritmetico}. Notiamo che, in entrambi i casi, il nuovo numero si può ottenere prendendo l'originale e dividendolo per $2^N$, dove $N$ è il valore dello shift. Questa è una proprietà degli shift di destra. Allo stesso modo, per gli shift di sinistra, basta moltiplicare il valore per $2^N$. Ad esempio, $-8_{10} = 111000 << 2 = 100000 = -32 = -8 \cdot 2^2$.  

\begin{Proposizione}[Calcolare numero decimale dopo uno shift]
\label{propos-caloclare-decimale-post-shift}
Preso un numero $A_{10}$, effettuando uno shift a \textbf{sinistra} di $N$ posizioni, troviamo che $\bm{A_{10} << N = A_{10} \cdot 2^n}$. Se invece si effettua uno shift a \textbf{destra}, troviamo che $\bm{A >> N = A / 2^N}$.  
\end{Proposizione}


\end{document}
